\paragraph{Source check for the fixed-second-coordinate assertion.}
Lochak--Schneps, Lemma~9 on p.~583 of
\cite{R-LS}, asserts that a fixed $f$ occurs only with the two
multipliers $\lambda$ and $-\lambda$.  The displayed class-two
identity $\lambda^2=24r+1$ determines the squares of the prime
coordinates.  The second proof in that source adds the required
full-group equation, whose elementary verification is as follows.

Comparing $(\lambda,f)$ and $(\mu,f)$ by the exact group law
gives $(\lambda^{-1}\mu,1)$.  Set
$n=(\lambda^{-1}\mu-1)/2\in\widehat\Z$.
The three-cycle relation is
\[
                       x^n(xy)^{-n}y^n=1.                 \tag{LS.1}
\]
In the class-two Heisenberg quotients its central coordinate
gives $n(n+1)=0$.  Thus in each integral domain $\Z_p$ one
has $n_p=0$ or $n_p=-1$.  These choices are linked by actual
finite quotients: for every odd prime $q$, evaluate (LS.1) in
\[
 D_q=\langle r,s\mid r^q=s^2=1,\ srs=r^{-1}\rangle,
                         \qquad x=s,\quad y=r.
\]
Both $x$ and $xy$ have order two.  If $n_2=0$, (LS.1) is
$r^{n_q}=1$, forcing $n_q=0$.  If $n_2=-1$, it is
$r^{1+n_q}=1$, forcing $n_q=-1$.  Hence $n$ is globally
zero or globally minus one, proving $\mu=\lambda$ or
$\mu=-\lambda$.  Both occur, because multiplication on the
right by $(-1,1)$ leaves the second coordinate unchanged.

This verifies the source's fixed-$f$ conclusion by exact
quotient morphisms.  Powers of a general torsion element
have second coordinates $f\alpha(f)\cdots\alpha^{j-1}(f)$;
the argument does not identify those coordinates with each
other.  D.32 supplies the separate argument restricting
finite subgroups, and leaves the stated involution questions
unchanged.
